% Options for packages loaded elsewhere
% Options for packages loaded elsewhere
\PassOptionsToPackage{unicode}{hyperref}
\PassOptionsToPackage{hyphens}{url}
\PassOptionsToPackage{dvipsnames,svgnames,x11names}{xcolor}
\PassOptionsToPackage{space}{xeCJK}
%
\documentclass[
  12pt,
  letterpaper,
  DIV=11,
  numbers=noendperiod]{scrartcl}
\usepackage{xcolor}
\usepackage{amsmath,amssymb}
\setcounter{secnumdepth}{-\maxdimen} % remove section numbering
\usepackage{iftex}
\ifPDFTeX
  \usepackage[T1]{fontenc}
  \usepackage[utf8]{inputenc}
  \usepackage{textcomp} % provide euro and other symbols
\else % if luatex or xetex
  \usepackage{unicode-math} % this also loads fontspec
  \defaultfontfeatures{Scale=MatchLowercase}
  \defaultfontfeatures[\rmfamily]{Ligatures=TeX,Scale=1}
\fi
\usepackage{lmodern}
\ifPDFTeX\else
  % xetex/luatex font selection
  \setsansfont[]{Century Schoolbook}
  \ifXeTeX
    \usepackage{xeCJK}
    \setCJKmainfont[]{Songti TC}
  \fi
  \ifLuaTeX
    \usepackage[]{luatexja-fontspec}
    \setmainjfont[]{Songti TC}
  \fi
\fi
% Use upquote if available, for straight quotes in verbatim environments
\IfFileExists{upquote.sty}{\usepackage{upquote}}{}
\IfFileExists{microtype.sty}{% use microtype if available
  \usepackage[]{microtype}
  \UseMicrotypeSet[protrusion]{basicmath} % disable protrusion for tt fonts
}{}
\makeatletter
\@ifundefined{KOMAClassName}{% if non-KOMA class
  \IfFileExists{parskip.sty}{%
    \usepackage{parskip}
  }{% else
    \setlength{\parindent}{0pt}
    \setlength{\parskip}{6pt plus 2pt minus 1pt}}
}{% if KOMA class
  \KOMAoptions{parskip=half}}
\makeatother
% Make \paragraph and \subparagraph free-standing
\makeatletter
\ifx\paragraph\undefined\else
  \let\oldparagraph\paragraph
  \renewcommand{\paragraph}{
    \@ifstar
      \xxxParagraphStar
      \xxxParagraphNoStar
  }
  \newcommand{\xxxParagraphStar}[1]{\oldparagraph*{#1}\mbox{}}
  \newcommand{\xxxParagraphNoStar}[1]{\oldparagraph{#1}\mbox{}}
\fi
\ifx\subparagraph\undefined\else
  \let\oldsubparagraph\subparagraph
  \renewcommand{\subparagraph}{
    \@ifstar
      \xxxSubParagraphStar
      \xxxSubParagraphNoStar
  }
  \newcommand{\xxxSubParagraphStar}[1]{\oldsubparagraph*{#1}\mbox{}}
  \newcommand{\xxxSubParagraphNoStar}[1]{\oldsubparagraph{#1}\mbox{}}
\fi
\makeatother


\usepackage{longtable,booktabs,array}
\newcounter{none} % for unnumbered tables
\usepackage{calc} % for calculating minipage widths
% Correct order of tables after \paragraph or \subparagraph
\usepackage{etoolbox}
\makeatletter
\patchcmd\longtable{\par}{\if@noskipsec\mbox{}\fi\par}{}{}
\makeatother
% Allow footnotes in longtable head/foot
\IfFileExists{footnotehyper.sty}{\usepackage{footnotehyper}}{\usepackage{footnote}}
\makesavenoteenv{longtable}
\usepackage{graphicx}
\makeatletter
\newsavebox\pandoc@box
\newcommand*\pandocbounded[1]{% scales image to fit in text height/width
  \sbox\pandoc@box{#1}%
  \Gscale@div\@tempa{\textheight}{\dimexpr\ht\pandoc@box+\dp\pandoc@box\relax}%
  \Gscale@div\@tempb{\linewidth}{\wd\pandoc@box}%
  \ifdim\@tempb\p@<\@tempa\p@\let\@tempa\@tempb\fi% select the smaller of both
  \ifdim\@tempa\p@<\p@\scalebox{\@tempa}{\usebox\pandoc@box}%
  \else\usebox{\pandoc@box}%
  \fi%
}
% Set default figure placement to htbp
\def\fps@figure{htbp}
\makeatother





\setlength{\emergencystretch}{3em} % prevent overfull lines

\providecommand{\tightlist}{%
  \setlength{\itemsep}{0pt}\setlength{\parskip}{0pt}}



 


\KOMAoption{captions}{tableheading}
\makeatletter
\@ifpackageloaded{caption}{}{\usepackage{caption}}
\AtBeginDocument{%
\ifdefined\contentsname
  \renewcommand*\contentsname{Table of contents}
\else
  \newcommand\contentsname{Table of contents}
\fi
\ifdefined\listfigurename
  \renewcommand*\listfigurename{List of Figures}
\else
  \newcommand\listfigurename{List of Figures}
\fi
\ifdefined\listtablename
  \renewcommand*\listtablename{List of Tables}
\else
  \newcommand\listtablename{List of Tables}
\fi
\ifdefined\figurename
  \renewcommand*\figurename{Figure}
\else
  \newcommand\figurename{Figure}
\fi
\ifdefined\tablename
  \renewcommand*\tablename{Table}
\else
  \newcommand\tablename{Table}
\fi
}
\@ifpackageloaded{float}{}{\usepackage{float}}
\floatstyle{ruled}
\@ifundefined{c@chapter}{\newfloat{codelisting}{h}{lop}}{\newfloat{codelisting}{h}{lop}[chapter]}
\floatname{codelisting}{Listing}
\newcommand*\listoflistings{\listof{codelisting}{List of Listings}}
\makeatother
\makeatletter
\makeatother
\makeatletter
\@ifpackageloaded{caption}{}{\usepackage{caption}}
\@ifpackageloaded{subcaption}{}{\usepackage{subcaption}}
\makeatother
\usepackage{bookmark}
\IfFileExists{xurl.sty}{\usepackage{xurl}}{} % add URL line breaks if available
\urlstyle{same}
\hypersetup{
  pdftitle={Institutional Investors: Theory and Evidence},
  pdfauthor={Calvin J. Chiou},
  colorlinks=true,
  linkcolor={blue},
  filecolor={Maroon},
  citecolor={Blue},
  urlcolor={Blue},
  pdfcreator={LaTeX via pandoc}}


\title{Institutional Investors: Theory and Evidence}
\usepackage{etoolbox}
\makeatletter
\providecommand{\subtitle}[1]{% add subtitle to \maketitle
  \apptocmd{\@title}{\par {\large #1 \par}}{}{}
}
\makeatother
\subtitle{Fall AY115 \textbar{} NCCU College of Commerce}
\author{}
\date{}
\begin{document}
\maketitle

\renewcommand*\contentsname{Table of contents}
{
\setcounter{tocdepth}{2}
\tableofcontents
}

\subsection{Course Information}\label{course-information}

\textbf{Course Code:} 357854001

\textbf{Time and Venue:} Thursday 16:10--18:00, Venue TBA

\textbf{Lecturer:} Asst. Prof.~Calvin J. Chiou (邱健嘉)

\textbf{Office:} Commerce Building 261233

\textbf{Office Hours:} By appointment

\textbf{Email:} \texttt{jjchiou@nccu.edu.tw}

\begin{center}\rule{0.5\linewidth}{0.5pt}\end{center}

\subsection{Course Description}\label{course-description}

This course explores the economic role and strategic behavior of
institutional investors in financial markets and corporate governance.
We cover different types of institutional investors (e.g., mutual funds,
pension funds, and insurance companies) examining their investment
objectives and constraints, and their influence on firm policies, market
efficiency, and ESG practices. Special attention is given to the Taiwan
and broader Asian institutional investor landscape to connect theory to
the local context.

A distinguishing feature of this course is its integration of empirical
methods with theory. Five sessions include instructor-led R labs in
which students implement fund performance measures, holdings-based
analyses, and factor model regressions directly in R using CRSP mutual
fund data and Kenneth French's public factor library. Lecture notes for
these sessions are distributed as R Quarto documents that students can
run and extend independently.

\subsection{Learning Objectives}\label{learning-objectives}

By the end of the course, students will be able to:

\begin{enumerate}
\def\labelenumi{\arabic{enumi}.}
\tightlist
\item
  Identify and distinguish the major types of institutional investors
  and their structural incentives.
\item
  Analyze how institutional investors affect corporate governance, firm
  behavior, and market prices.
\item
  Critically evaluate empirical research on institutional investor
  strategies, performance, and market impact.
\item
  Recognize the key identification challenges in this literature and
  discuss how leading papers address them.
\item
  Implement core fund performance and holdings-based measures in R,
  including factor model alphas, DGTW benchmarks, active share, churn
  ratio, and window dressing indices.
\item
  Develop and present an original research idea or critical literature
  review in the institutional investor domain.
\end{enumerate}

\subsection{Course Materials}\label{course-materials}

\textbf{Academic papers:} Readings are drawn from published and working
papers across top finance journals including the \emph{Journal of
Finance} (JF), \emph{Journal of Financial Economics} (JFE), \emph{Review
of Financial Studies} (RFS), \emph{Management Science} (MS),
\emph{Journal of Financial and Quantitative Analysis} (JFQA), and
\emph{Review of Finance} (RF). A complete reading list is provided at
the end of this syllabus.

\textbf{R lecture notes:} Empirical lab sessions are accompanied by
\texttt{.qmd} lecture notes distributed via Moodle. Students should have
R (≥ 4.3) and RStudio (or Positron) installed.

\subsection{Course Format}\label{course-format}

This is a graduate research seminar combining three modes of
instruction:

\begin{itemize}
\tightlist
\item
  \textbf{Lecture and discussion} (\textasciitilde50\%): The instructor
  introduces the week's topic, reviews key papers, and frames open
  research questions.
\item
  \textbf{Student presentations} (\textasciitilde30\%): Students present
  assigned papers from the asset management literature.
\item
  \textbf{Empirical R labs} (\textasciitilde20\%): Five instructor-led
  sessions demonstrating how to implement measures from the readings in
  R. Lab notes are distributed on the instructor's website.
\end{itemize}

Active participation is expected. Students are encouraged to read at
least the abstract, introduction, and conclusion of all assigned papers
before class.

\subsection{Software and Computing}\label{software-and-computing}

All empirical work in this course uses \textbf{R} via \textbf{RStudio}
or \textbf{Positron}. Lecture notes are written in \textbf{R Quarto} and
are self-contained: they load data, run all analyses, and produce
formatted output. Labs 1--3 require no proprietary data access. Labs
4--5 use CRSP Mutual Fund Holdings; data files will be publicly
available or distributed by the instructor via Moodle.

Students are expected to have R and the required packages installed
before Week 2.

\subsection{Assessment}\label{assessment}

{\def\LTcaptype{none} % do not increment counter
\begin{longtable}[]{@{}ll@{}}
\toprule\noalign{}
Component & Weight \\
\midrule\noalign{}
\endhead
\bottomrule\noalign{}
\endlastfoot
Class Participation \& Discussion & 20\% \\
Paper Presentations & 30\% \\
Final Written Report & 20\% \\
Final Presentation & 30\% \\
\end{longtable}
}

\subsubsection{Paper Presentations}\label{paper-presentations}

Each student will present \textbf{two papers} during the semester (or
one extended presentation covering a topic area). Presentation slots are
allocated during Week 1. Each presentation should:

\begin{itemize}
\tightlist
\item
  Summarize the research question, data, methodology, and main findings.
\item
  Situate the paper within the broader literature covered in the course.
\item
  Raise at least two discussion questions for the class.
\item
  Be approximately \textbf{20--25 minutes}, followed by 10 minutes of
  Q\&A.
\end{itemize}

Presentations may be individual or in pairs (to be confirmed based on
enrollment). Slides are required.

\subsubsection{Final Report and
Presentation}\label{final-report-and-presentation}

\textbf{Each student} completes a final report demonstrating their
ability to engage with the institutional investor literature. The report
should focus on a specific topic and take one of the following forms:

\begin{itemize}
\tightlist
\item
  A \textbf{literature review and synthesis} of a focused sub-topic
  (minimum 10 papers).
\item
  A \textbf{preliminary} \textbf{research project} with a clearly stated
  question, data strategy, expected contribution, and empirical results.
\item
  A \textbf{replication and extension} of an existing empirical paper,
  optionally using the R methods introduced in the lab sessions.
\end{itemize}

\textbf{Key deadlines:}

{\def\LTcaptype{none} % do not increment counter
\begin{longtable}[]{@{}
  >{\raggedright\arraybackslash}p{(\linewidth - 2\tabcolsep) * \real{0.7660}}
  >{\raggedright\arraybackslash}p{(\linewidth - 2\tabcolsep) * \real{0.2340}}@{}}
\toprule\noalign{}
\begin{minipage}[b]{\linewidth}\raggedright
Milestone
\end{minipage} & \begin{minipage}[b]{\linewidth}\raggedright
Deadline
\end{minipage} \\
\midrule\noalign{}
\endhead
\bottomrule\noalign{}
\endlastfoot
Topic proposal (1 paragraph via Moodle or email) & End of Week 8 \\
Final written report (≤ 10 double-spaced pages, excl. references) & End
of Week 17 \\
Final in-class presentation (20--30 min) & Weeks 15--16 \\
\end{longtable}
}

\begin{center}\rule{0.5\linewidth}{0.5pt}\end{center}

\subsection{Course Schedule}\label{course-schedule}

Weeks marked \textbf{{[}Lab{]}} include an instructor-led R session in
the second half of class. Lab notes (\texttt{.qmd}) are distributed on
Moodle before each lab week. Theory lecture and paper discussion occupy
the first half of the session.

\begin{longtable}[]{@{}
  >{\raggedright\arraybackslash}p{(\linewidth - 8\tabcolsep) * \real{0.0532}}
  >{\raggedright\arraybackslash}p{(\linewidth - 8\tabcolsep) * \real{0.1064}}
  >{\raggedright\arraybackslash}p{(\linewidth - 8\tabcolsep) * \real{0.2872}}
  >{\raggedright\arraybackslash}p{(\linewidth - 8\tabcolsep) * \real{0.3723}}
  >{\raggedright\arraybackslash}p{(\linewidth - 8\tabcolsep) * \real{0.1809}}@{}}
\caption{Schedule is tentative and subject to change. Lab sessions
assume basic familiarity with R and
\protect\texttt{tidyverse}.}\tabularnewline
\toprule\noalign{}
\begin{minipage}[b]{\linewidth}\raggedright
Week
\end{minipage} & \begin{minipage}[b]{\linewidth}\raggedright
Type
\end{minipage} & \begin{minipage}[b]{\linewidth}\raggedright
Topic
\end{minipage} & \begin{minipage}[b]{\linewidth}\raggedright
Key Readings
\end{minipage} & \begin{minipage}[b]{\linewidth}\raggedright
Notes
\end{minipage} \\
\midrule\noalign{}
\endfirsthead
\toprule\noalign{}
\begin{minipage}[b]{\linewidth}\raggedright
Week
\end{minipage} & \begin{minipage}[b]{\linewidth}\raggedright
Type
\end{minipage} & \begin{minipage}[b]{\linewidth}\raggedright
Topic
\end{minipage} & \begin{minipage}[b]{\linewidth}\raggedright
Key Readings
\end{minipage} & \begin{minipage}[b]{\linewidth}\raggedright
Notes
\end{minipage} \\
\midrule\noalign{}
\endhead
\bottomrule\noalign{}
\endlastfoot
1 & Lecture & Course Overview and Motivation & Bebchuk, Cohen \& Hirst
(2017) & Presentation slots assigned; R setup check \\
2 & Lecture & Types of Institutional Investors and Asset Concentration
Trends & Gompers \& Metrick (2001); Ferreira \& Matos (2008) &
Taiwan/Asian context \\
3 & \textbf{{[}Lab 1{]}} & Measuring Fund Performance I: Carhart
Four-Factor Model & Carhart (1997); Wermers (2011) & \emph{Lab: factor
model estimation} \\
4 & Lecture & Fund Performance and the Luck vs.~Skill Debate & Fama \&
French (2010); Berk \& van Binsbergen (2015) & Student presentation \\
5 & \textbf{{[}Lab 2{]}} & Measuring Fund Performance II: DGTW
Characteristic-Adjusted Returns & Daniel et al.~(1997); Wermers (2000) &
\emph{Lab: constructing size × BM × momentum benchmarks} \\
6 & Lecture & Benchmarking, Style Drift, and Active Share & Bollen \&
Busse (2005); Cremers \& Petajisto (2009) & Student presentation \\
7 & Lecture & Tournament Behavior, Risk Shifting, and Fund Flows & Brown
et al.~(1996); Sirri \& Tufano (1998) & Student presentation \\
8 & Lecture & Manager Incentives and Fund Family Conflicts & Gaspar et
al.~(2006); Bhattacharya et al.~(2013) & \textbf{Topic proposals due};
student presentation \\
9 & \textbf{{[}Lab 3{]}} & Measuring Fund Performance III: Bootstrap
Inference for Luck vs.~Skill & Fama \& French (2010) & \emph{Lab: joint
block bootstrap; actual vs.~simulated alpha distribution in R} \\
10 & Lecture & Fund Holdings and Trading Behavior & Kacperczyk, Sialm \&
Zheng (2005); Wermers (1999) & Student presentation \\
11 & \textbf{{[}Lab 4{]}} & Holdings-Based Measures: Active Share, Churn
Ratio, and Window Dressing & Cremers \& Petajisto (2009); Agarwal, Gay
\& Ling (2014) & \emph{Lab: computing active share, churn ratio, and
window dressing} \\
12 & Lecture & Career Concerns, Performance Manipulation, and Behavioral
Biases & Kempf et al.~(2009); Frazzini \& Lamont (2008) & Student
presentation \\
13 & Lecture & Institutional Monitoring, Voting, and Proxy Engagement &
Gillan \& Starks (2000); Iliev \& Lowry (2015) & Student presentation \\
14 & \textbf{{[}Lab 5{]}} & Holdings-Based Analysis: Institutional
Ownership and Common Ownership Measures & Appel et al.~(2016); Azar et
al.~(2018) & \emph{Lab: ownership concentration, MHHI common
ownership} \\
15 & Lecture & ESG Investing, Climate Risk, and Long-Term Stewardship &
Chen, Dong \& Lin (2020); Starks, Venkat \& Zhu (2025) & Student
presentation \\
16 & Presentations & Final Presentations & --- & \textbf{Written reports
due end of week} \\
\end{longtable}

\subsection{R Lab Sessions: Detailed
Outlines}\label{r-lab-sessions-detailed-outlines}

\subsubsection{Lab 1 --- Carhart Four-Factor Model (Week
3)}\label{lab-1-carhart-four-factor-model-week-3}

\textbf{Objective:} Estimate factor model alphas and factor loadings for
a cross-section of ETFs using publicly available data.

\textbf{Data:} \texttt{tidyquant} (ETF prices); \texttt{frenchdata}
(Fama-French 5 factors + momentum)

\textbf{Skills covered:}

\begin{itemize}
\tightlist
\item
  Downloading adjusted price data with \texttt{tidyquant}
\item
  Obtaining Fama-French and momentum factors with \texttt{frenchdata}
\item
  Computing monthly excess returns and aligning on year-month
\item
  Estimating CAPM, FF3, and Carhart 4-factor regressions with
  \texttt{broom::tidy()}
\item
  Interpreting alpha, factor loadings, and R² across model
  specifications
\item
  Visualizing factor loadings and alpha comparisons with
  \texttt{ggplot2}
\end{itemize}

\textbf{Reference:} Lecture Note 1; Carhart (1997); Wermers (2011,
Section 3)

\subsubsection{Lab 2 --- DGTW Characteristic-Adjusted Returns (Week
5)}\label{lab-2-dgtw-characteristic-adjusted-returns-week-5}

\textbf{Objective:} Implement the Daniel, Grinblatt, Titman \& Wermers
(1997) characteristic-adjusted return benchmark and decompose fund
performance into stock-selection and style components.

\textbf{Data:} \texttt{tidyquant} (stock prices and characteristics);
\texttt{frenchdata} (size/BM breakpoints)

\textbf{Skills covered:}

\begin{itemize}
\tightlist
\item
  Constructing size × book-to-market × momentum benchmark portfolios
  from public data
\item
  Computing the DGTW Characteristic Selectivity (CS) and Characteristic
  Timing (CT) measures
\item
  Calculating fund-level flow and turnover ratio from return and TNA
  data
\item
  Approximating Active Share from ETF holdings available via
  \texttt{tidyquant}
\end{itemize}

\textbf{Reference:} Lecture Note 2; Daniel et al.~(1997); Wermers (2000)

\subsubsection{Lab 3 --- Bootstrap Inference for Luck vs.~Skill (Week
9)}\label{lab-3-bootstrap-inference-for-luck-vs.-skill-week-9}

\textbf{Objective:} Replicate the bootstrap simulation logic of Fama \&
French (2010) to test whether the cross-sectional alpha distribution is
consistent with a zero-skill null hypothesis.

\textbf{Data:} CRSP mutual fund monthly net returns (or ETF proxies from
Lab 1)

\textbf{Skills covered:}

\begin{itemize}
\tightlist
\item
  Estimating Carhart alphas for a panel of funds using
  \texttt{group\_by()} + \texttt{broom::tidy()}
\item
  Implementing joint monthly block bootstrap resampling (preserving
  cross-fund correlation structure)
\item
  Comparing actual vs.~simulated alpha percentiles across the
  distribution
\item
  Visualizing the luck vs.~skill decomposition with density plots and
  percentile bands
\end{itemize}

\textbf{Reference:} Lecture Note 3; Fama \& French (2010, Section III
and Appendix)

\subsubsection{Lab 4 --- Holdings-Based Measures: Active Share, Churn
Ratio, and Window Dressing (Week
11)}\label{lab-4-holdings-based-measures-active-share-churn-ratio-and-window-dressing-week-11}

\textbf{Objective:} Use CRSP Mutual Fund Holdings to compute three
widely-used measures of fund trading behavior and portfolio activeness.

\textbf{Data:} CRSP Mutual Fund Holdings (quarterly stock-level
positions); CRSP monthly stock returns

\textbf{Skills covered:}

\begin{itemize}
\tightlist
\item
  Loading and cleaning CRSP mutual fund holdings data in R
\item
  Computing \textbf{Active Share} (Cremers \& Petajisto 2009): fraction
  of portfolio holdings that differ from a benchmark index
\item
  Computing the \textbf{Churn Ratio} (Gaspar, Massa \& Matos 2005): a
  holdings-based measure of portfolio turnover that captures short-term
  trading intensity
\item
  Constructing a \textbf{Window Dressing Index} (Agarwal, Gay \& Ling
  2014): detecting quarter-end inflation of recent winners and dumping
  of recent losers in disclosed holdings
\item
  Merging holdings data with CRSP return data using \texttt{wfundno} /
  \texttt{permno} links
\item
  Summarizing and visualizing cross-fund distributions of each measure
  over time
\end{itemize}

\textbf{Reference:} Lecture Note 4; Cremers \& Petajisto (2009);
Agarwal, Gay \& Ling (2014); Gaspar, Massa \& Matos (2006)

\subsubsection{Lab 5 --- Institutional Ownership and Common Ownership
Measures (Week
14)}\label{lab-5-institutional-ownership-and-common-ownership-measures-week-14}

\textbf{Objective:} Construct firm-level institutional ownership and
common ownership measures from holdings data, and run panel regressions
relating ownership to firm outcomes.

\textbf{Data:} CRSP Mutual Fund Holdings (or 13F proxy); CRSP/Compustat
merged panel

\textbf{Skills covered:}

\begin{itemize}
\tightlist
\item
  Aggregating fund-level holdings to firm-level institutional ownership
  (IO) percentage
\item
  Computing ownership concentration: Herfindahl-Hirschman Index (HHI) of
  fund ownership per firm
\item
  Constructing the modified HHI (MHHI delta) of common ownership
  following Azar et al.~(2018)
\item
  Merging ownership data with stock return and accounting data for panel
  analysis
\item
  Running firm-level panel regressions with two-way (firm + year) fixed
  effects using \texttt{fixest::feols()}
\item
  Interpreting coefficient magnitudes and standard error clustering
  choices
\end{itemize}

\textbf{Reference:} Lecture Note 5; Appel et al.~(2016); Azar et
al.~(2018)

\subsection{Reading List}\label{reading-list}

\subsubsection{I. Overview and Types of Institutional
Investors}\label{i.-overview-and-types-of-institutional-investors}

Bebchuk, L. A., Cohen, A., \& Hirst, S. (2017). The agency problems of
institutional investors. \emph{Journal of Economic Perspectives},
\emph{31}(3), 89--112.

Bushee, B. J. (1998). The influence of institutional investors on myopic
R\&D investment behavior. \emph{The Accounting Review}, \emph{73}(3),
305--333.

Bushee, B. J. (2001). Do institutional investors prefer near-term
earnings over long-run value? \emph{Contemporary Accounting Research},
\emph{18}(2), 207--246.

Ferreira, M. A., \& Matos, P. (2008). The colors of investors' money:
The role of institutional investors around the world. \emph{Journal of
Financial Economics}, \emph{88}(3), 499--533.

Gompers, P. A., \& Metrick, A. (2001). Institutional investors and
equity prices. \emph{Quarterly Journal of Economics}, \emph{116}(1),
229--259.

Huang, R. D., \& Shiu, C. Y. (2009). Local effects of foreign ownership
in an emerging financial market: Evidence from qualified foreign
institutional investors in Taiwan. \emph{Financial Management},
\emph{38}(3), 567-602. \emph{(Taiwan context)}

\subsubsection{II. Fund Performance Measurement and
Evaluation}\label{ii.-fund-performance-measurement-and-evaluation}

Berk, J. B., \& Van Binsbergen, J. H. (2015). Measuring skill in the
mutual fund industry. \emph{Journal of Financial Economics},
\emph{118}(1), 1--20.

Bollen, N. P., \& Busse, J. A. (2005). Short-term persistence in mutual
fund performance. \emph{Review of Financial Studies}, \emph{18}(2),
569--597.

Carhart, M. M. (1997). On persistence in mutual fund performance.
\emph{Journal of Finance}, \emph{52}(1), 57--82.

Cremers, K. M., \& Petajisto, A. (2009). How active is your fund
manager? A new measure that predicts performance. \emph{Review of
Financial Studies}, \emph{22}(9), 3329--3365.

Daniel, K., Grinblatt, M., Titman, S., \& Wermers, R. (1997). Measuring
mutual fund performance with characteristic-based benchmarks.
\emph{Journal of Finance}, \emph{52}(3), 1035--1058.

Fama, E. F., \& French, K. R. (2010). Luck versus skill in the
cross-section of mutual fund returns. \emph{Journal of Finance},
\emph{65}(5), 1915--1947.

Wermers, R. (2000). Mutual fund performance: An empirical decomposition
into stock-picking talent, style, transactions costs, and expenses.
\emph{Journal of Finance}, \emph{55}(4), 1655--1703.

Wermers, R. (2011). Performance measurement of mutual funds, hedge
funds, and institutional accounts. \emph{Annual Review of Financial
Economics}, \emph{3}, 537--574.

\subsubsection{III. Fund Holdings, Trading Behavior, and
Activeness}\label{iii.-fund-holdings-trading-behavior-and-activeness}

Agarwal, V., Gay, G. D., \& Ling, L. (2014). Window dressing in mutual
funds. \emph{Review of Financial Studies}, \emph{27}(11), 3133--3170.

Lou, D. (2012). A flow-based explanation for return predictability.
\emph{Review of Financial Studies}, \emph{25}(12), 3457-3489.

Gaspar, J. M., Massa, M., \& Matos, P. (2005). Shareholder investment
horizons and the market for corporate control. \emph{Journal of
Financial Economics}, \emph{76}(1), 135--165. \emph{(Churn ratio
methodology)}

Kacperczyk, M., Sialm, C., \& Zheng, L. (2005). On the industry
concentration of actively managed equity mutual funds. \emph{Journal of
Finance}, \emph{60}(4), 1983--2011.

Kacperczyk, M., Sialm, C., \& Zheng, L. (2008). Unobserved actions of
mutual funds. \emph{Review of Financial Studies}, \emph{21}(6),
2379--2416.

Wermers, R. (1999). Mutual fund herding and the impact on stock prices.
\emph{Journal of Finance}, \emph{54}(2), 581--622.

\subsubsection{IV. Agency Problems and Incentives in Fund
Management}\label{iv.-agency-problems-and-incentives-in-fund-management}

Bhattacharya, U., Lee, J. H., \& Pool, V. K. (2013). Conflicting family
values in mutual fund families. \emph{Journal of Finance}, \emph{68}(1),
173--200.

Brown, K. C., Harlow, W. V., \& Starks, L. T. (1996). Of tournaments and
temptations: An analysis of managerial incentives in the mutual fund
industry. \emph{Journal of Finance}, \emph{51}(1), 85--110.

Gaspar, J. M., Massa, M., \& Matos, P. (2006). Favoritism in mutual fund
families? Evidence on strategic cross-fund subsidization. \emph{Journal
of Finance}, \emph{61}(1), 73--104.

Kempf, A., Ruenzi, S., \& Thiele, T. (2009). Employment risk,
compensation incentives, and managerial risk taking: Evidence from the
mutual fund industry. \emph{Journal of Financial Economics},
\emph{92}(1), 92--108.

Sirri, E. R., \& Tufano, P. (1998). Costly search and mutual fund flows.
\emph{Journal of Finance}, \emph{53}(5), 1589--1622.

\subsubsection{V. Behavioral Biases of
Investors}\label{v.-behavioral-biases-of-investors}

Agarwal, V., Jiang, L., \& Wen, Q. (2022). Why do mutual funds hold
lottery stocks?. \emph{Journal of Financial and Quantitative Analysis},
\emph{57}(3), 825-856.

Barber, B. M., Odean, T., \& Zheng, L. (2005). Out of sight, out of
mind: The effects of expenses on mutual fund flows. \emph{Journal of
Business}, \emph{78}(6), 2095--2119.

Ben-David, I., Franzoni, F., Kim, B., \& Moussawi, R. (2021).
Competition for attention in the ETF space. \emph{Review of Financial
Studies}, \emph{35}(9), 4382--4427.

Frazzini, A., \& Lamont, O. A. (2008). Dumb money: Mutual fund flows and
the cross-section of stock returns. \emph{Journal of Financial
Economics}, \emph{88}(2), 299--322.

\subsubsection{VI. Institutional Monitoring, Voting, and
Engagement}\label{vi.-institutional-monitoring-voting-and-engagement}

Chen, X., Harford, J., \& Li, K. (2007). Monitoring: Which institutions
matter? \emph{Journal of Financial Economics}, \emph{86}(2), 279--305.

Chung, K. H., \& Zhang, H. (2011). Corporate governance and
institutional ownership. \emph{Journal of Financial and Quantitative
Analysis}, \emph{46}(1), 247--273.

Gillan, S. L., \& Starks, L. T. (2000). Corporate governance proposals
and shareholder activism: The role of institutional investors.
\emph{Journal of Financial Economics}, \emph{57}(2), 275--305.

Iliev, P., \& Lowry, M. (2015). Are mutual funds active voters?
\emph{Review of Financial Studies}, \emph{28}(2), 446--485.

Matvos, G., \& Ostrovsky, M. (2010). Heterogeneity and peer effects in
mutual fund proxy voting. \emph{Journal of Financial Economics},
\emph{98}(1), 90--112.

\subsubsection{VII. Passive Investing and Common
Ownership}\label{vii.-passive-investing-and-common-ownership}

Appel, I. R., Gormley, T. A., \& Keim, D. B. (2016). Passive investors,
not passive owners. \emph{Journal of Financial Economics},
\emph{121}(1), 111--141.

Azar, J., Schmalz, M. C., \& Tecu, I. (2018). Anticompetitive effects of
common ownership. \emph{Journal of Finance}, \emph{73}(4), 1513--1565.

Bebchuk, L., \& Hirst, S. (2019). The specter of the giant three.
\emph{Boston University Law Review}, \emph{99}(3), 721--741.

He, J., \& Huang, J. (2017). Product market competition in a world of
cross-ownership: Evidence from institutional blockholdings. \emph{Review
of Financial Studies}, \emph{30}(8), 2674--2718.

Heath, D., Macciocchi, D., Michaely, R., \& Ringgenberg, M. C. (2022).
Do index funds monitor? \emph{Review of Financial Studies},
\emph{35}(1), 91--131.

Lewellen, J., \& Lowry, M. (2021). Does common ownership really increase
firm coordination? \emph{Journal of Financial Economics}, \emph{141}(1),
322--344.

Massa, M., Schumacher, D., \& Wang, Y. (2021). Who is afraid of
BlackRock? \emph{Review of Financial Studies}, \emph{34}(4), 1987--2044.

\subsubsection{VIII. ESG Investing, Climate Risk, and Long-Term
Stewardship}\label{viii.-esg-investing-climate-risk-and-long-term-stewardship}

Chen, T., Dong, H., \& Lin, C. (2020). Institutional shareholders and
corporate social responsibility. \emph{Journal of Financial Economics},
\emph{135}(2), 483--504.

Dyck, A., Lins, K. V., Roth, L., \& Wagner, H. F. (2019). Do
institutional investors drive corporate social responsibility?
International evidence. \emph{Journal of Financial Economics},
\emph{131}(3), 693--714.

Hwang, C. Y., Titman, S., \& Wang, Y. (2022). Investor tastes, corporate
behavior, and stock returns: An analysis of corporate social
responsibility. \emph{Management Science}, \emph{68}(10), 7131--7152.

Kim, S., \& Yoon, A. (2023). Analyzing active fund managers' commitment
to ESG: Evidence from the United Nations Principles for Responsible
Investment. \emph{Management Science}, \emph{69}(2), 741--758.

Krueger, P., Sautner, Z., \& Starks, L. T. (2020). The importance of
climate risks for institutional investors. \emph{Review of Financial
Studies}, \emph{33}(3), 1067--1111.

Starks, L. T., Venkat, P., \& Zhu (2025). Corporate ESG profiles and
investor horizons. \emph{Journal of Finance}, \emph{81}(2), 603-642.




\end{document}
